\documentclass[10pt, letterpaper]{article}

% Packages:
\usepackage[
    ignoreheadfoot, % set margins without considering header and footer
    top=0.5 cm, % seperation between body and page edge from the top
    bottom=0.5 cm, % seperation between body and page edge from the bottom
    left=1 cm, % seperation between body and page edge from the left
    right=1 cm, % seperation between body and page edge from the right
    footskip=1.0 cm, % seperation between body and footer
    % showframe % for debugging
]{geometry} % for adjusting page geometry
\usepackage{titlesec} % for customizing section titles
\usepackage{tabularx} % for making tables with fixed width columns
\usepackage{array} % tabularx requires this
\usepackage[dvipsnames]{xcolor} % for coloring text
\definecolor{primaryColor}{RGB}{0, 0, 0} % define primary color
\usepackage{enumitem} % for customizing lists
\usepackage{fontawesome5} % for using icons
\usepackage{amsmath} % for math
\usepackage[
    pdftitle={Tzuching's CV},
    pdfauthor={Tzuching Liao},
    pdfcreator={LaTeX with RenderCV},
    colorlinks=true,
    urlcolor=primaryColor
]{hyperref} % for links, metadata and bookmarks
\usepackage[pscoord]{eso-pic} % for floating text on the page
\usepackage{calc} % for calculating lengths
\usepackage{bookmark} % for bookmarks
\usepackage{lastpage} % for getting the total number of pages
\usepackage{changepage} % for one column entries (adjustwidth environment)
\usepackage{paracol} % for two and three column entries
\usepackage{ifthen} % for conditional statements
\usepackage{needspace} % for avoiding page brake right after the section title
\usepackage{iftex} % check if engine is pdflatex, xetex or luatex




% Ensure that generate pdf is machine readable/ATS parsable:
\ifPDFTeX
    \input{glyphtounicode}
    \pdfgentounicode=1
    \usepackage[T1]{fontenc}
    \usepackage[utf8]{inputenc}
    \usepackage{lmodern}
\fi

\usepackage{charter}

% Some settings:
\raggedright
\AtBeginEnvironment{adjustwidth}{\partopsep0pt} % remove space before adjustwidth environment
\pagestyle{empty} % no header or footer
\setcounter{secnumdepth}{0} % no section numbering
\setlength{\parindent}{0pt} % no indentation
\setlength{\topskip}{0pt} % no top skip
\setlength{\columnsep}{0.15cm} % set column seperation
\pagenumbering{gobble} % no page numbering

\titleformat{\section}{\needspace{4\baselineskip}\bfseries\large}{}{0pt}{}[\vspace{1pt}\titlerule]

\titlespacing{\section}{
    % left space:
    -1pt
}{
    % top space:
    0.3 cm
}{
    % bottom space:
    0.2 cm
} % section title spacing

\renewcommand\labelitemi{$\vcenter{\hbox{\small$\bullet$}}$} % custom bullet points
\newenvironment{highlights}{
    \begin{itemize}[
        topsep=0.10 cm,
        parsep=0.10 cm,
        partopsep=0pt,
        itemsep=0pt,
        leftmargin=0 cm + 10pt
    ]
}{
    \end{itemize}
} % new environment for highlights


\newenvironment{highlightsforbulletentries}{
    \begin{itemize}[
        topsep=0.1 cm,
        parsep=0.1 cm,
        partopsep=0pt,
        itemsep=0pt,
        leftmargin=10pt
    ]
}{
    \end{itemize}
} % new environment for highlights for bullet entries

\newenvironment{onecolentry}{
    \begin{adjustwidth}{
        0 cm + 0.00001 cm
    }{
        0 cm + 0.00001 cm
    }
}{
    \end{adjustwidth}
} % new environment for one column entries

\newenvironment{twocolentry}[2][]{
    \onecolentry
    \def\secondColumn{#2}
    \setcolumnwidth{\fill, 4.5 cm}
    \begin{paracol}{2}
}{
    \switchcolumn \raggedleft \secondColumn
    \end{paracol}
    \endonecolentry
} % new environment for two column entries

\newenvironment{threecolentry}[3][]{
    \onecolentry
    \def\thirdColumn{#3}
    \setcolumnwidth{, \fill, 4.5 cm}
    \begin{paracol}{3}
    {\raggedright #2} \switchcolumn
}{
    \switchcolumn \raggedleft \thirdColumn
    \end{paracol}
    \endonecolentry
} % new environment for three column entries

\newenvironment{header}{
    \setlength{\topsep}{0pt}\par\kern\topsep\centering\linespread{1.5}
}{
    \par\kern\topsep
} % new environment for the header

\newcommand{\placelastupdatedtext}{% \placetextbox{<horizontal pos>}{<vertical pos>}{<stuff>}
  \AddToShipoutPictureFG*{% Add <stuff> to current page foreground
    \put(
        \LenToUnit{\paperwidth-2 cm-0 cm+0.05cm},
        \LenToUnit{\paperheight-1.0 cm}
    ){\vtop{{\null}\makebox[0pt][c]{
        \small\color{gray}\textit{Last updated in September 2024}\hspace{\widthof{Last updated in September 2024}}
    }}}%
  }%
}%

% save the original href command in a new command:
\let\hrefWithoutArrow\href

% new command for external links:


\begin{document}
    \newcommand{\AND}{\unskip
        \cleaders\copy\ANDbox\hskip\wd\ANDbox
        \ignorespaces
    }
    \newsavebox\ANDbox
    \sbox\ANDbox{$|$}

    \begin{header}
        \fontsize{25 pt}{25 pt}\selectfont Tzu-Ching Liao

        \vspace{5 pt}

        \normalsize
        \mbox{Los Angeles, United States}%
        \kern 5.0 pt%
        \AND%
        \kern 5.0 pt%
        \mbox{\href{mailto:angus3938@gmail.com}{angus3938@gmail.com}}%
        \kern 5.0 pt%
        \AND%
        \kern 5.0 pt%
        \mbox{\hrefWithoutArrow{tel:(213) 298-2245}{(213) 298-2245}}%
        \kern 5.0 pt%
        \AND%
        \kern 5.0 pt%
        \mbox{\href{https://www.linkedin.com/in/tzuching-liao-872516294/}{Linkedin}}%
        \kern 5.0 pt%
        \AND%
        \kern 5.0 pt%
        \mbox{\href{https://github.com/tliao730}{Github}}%
        \kern 5.0 pt%
        \AND%
    \end{header}

    \vspace{5 pt - 0.3 cm}

%-------------------------------------------------------------
    \section{Summary}
    Data Science MS candidate at USC designing LLM--orchestrated multi-model systems that route time series tasks across domains(healthcare, finance, transportation) to specialized forecasting modes.
%-------------------------------------------------------------
    \section{Education}

        \begin{twocolentry}{
            Aug. 2025 – May. 2027
        }
            \textbf{University of Southern California}, Master of Science in Applied Data Science \end{twocolentry}

        \vspace{0.10 cm}


        \begin{twocolentry}{
            Sep. 2020 – Jun. 2024
        }
            \textbf{National Taipei University}, Bachelor of Business Administration in Statistics  \end{twocolentry}
%------------------------------------------------------------

    \section{Professional Experience}

    \begin{onecolentry}
        \begin{twocolentry}{
            Mar. 2026 -
        }
        \textbf{\underline{Research Intern}} \\
        \textbf{Machine Learning and Data Mining Lab, Department of Computer Science, USC}
        \end{twocolentry}

        \begin{itemize}

            \item Designed and developed an LLM-as-orchestrator framework that routes natural language task descriptions to specialized time series foundation models, applying Mixture-of-Experts-inspired dispatch logic across healthcare, transportation, finance, and weather forecasting domains.
            \item Contributed to the TSOrchestra pipeline, integrating time series foundation models (e.g., Toto2) into a modular AI agent system that autonomously selects and coordinates forecasting models based on data characteristics and task requirements.
            \item Built and evaluated a traffic forecasting benchmark to validate the orchestration framework, demonstrating that LLM-driven model selection improves predictive performance over single-model baselines on real-world spatiotemporal data.

        \end{itemize}

    \end{onecolentry}


        \vspace{0.10 cm}
        \begin{onecolentry}
            \textbf{\underline{Data Engineer Intern}} \\
            \textbf{SYSTEX} - specializing in cloud services, AI, and integrated IT ecosystems, Taipei, Taiwan
            \begin{itemize}
                \item Architected a \textbf{high-velocity SQL Lakehouse} on \textbf{AWS} using Glue and PySpark , integrating automated schema validation and \textbf{CI/CD} pipelines for S3 infrastructure management to ensure data integrity. This optimization slashed data-to-analysis \textbf{latency} by \textbf{70\%} , enabling the business to shift from reactive reporting to real-time, proactive decision-making.
                \item Boosted executive decision efficiency by \textbf{50\% } by architecting low-latency data pipelines for live Tableau dashboards, replacing static reporting with real-time interactive analytics to enable faster strategic pivots.

                \end{itemize}
        \end{onecolentry}


    \begin{onecolentry}
        \begin{twocolentry}{
            Jul. 2022 – Jun. 2023 \\
            \underline{\textbf{Advisor: Prof. Hong-Long Wang}}
        }
        \textbf{\underline{Undergraduate Research Assistant}} \\
        \textbf{Department of Statistics, National Taipei University}
        \end{twocolentry}

        \begin{itemize}

            \item Accelerated retail strategy optimization in a growing market by engineering a high-scale predictive framework \textbf{(XGBoost/MLP)} on \textbf{3.5 billion records} , achieving \textbf{86\%} Recall for \textbf{repurchase forecasting}.
            \item By synthesizing \textbf{Association Analysis} and \textbf{K-Means} clustering, I transformed behavioral patterns into a 4-tier Customer Pyramid, enabling retailers to rapidly deploy personalized \textbf{CRM strategies} that maximize \textbf{customer retention}.


        \end{itemize}


    \end{onecolentry}

%-------------------------------------------------------------


    \section{Project Experience}
    % ==========================================================
         \begin{twocolentry}{

        }
            \textbf{Time Series Orchestra}\end{twocolentry}

        \vspace{0.10 cm}
        \begin{onecolentry}
            \begin{itemize}
                \item Built an \textbf{ensemble of state-of-the-art time series foundation models} (Moirai, Sundial, Toto, Toto 2.0, TimesFM) with an \textbf{LLM-based commander} conceived to coordinate model weighting, using cross-validation and \textbf{SLSQP} optimization to produce probabilistic forecasts. Ranked top \textbf{2 on MASE} on GIFT-Eval, a large-scale forecasting benchmark.
                \item Designed a model abstraction layer separating forecasting logic from evaluation interfacing, via a common Forecaster interface and a Predictor adapter layer, allowing every foundation model and the ensemble itself to be evaluated interchangeably against the same benchmarking harness — improving extensibility for adding or swapping models.
            \end{itemize}
        \end{onecolentry}
        % ==========================================================
        \begin{twocolentry}{

        }
            \textbf{Traffic Forecasting Benchmark}\end{twocolentry}

        \vspace{0.10 cm}
        \begin{onecolentry}
            \begin{itemize}
                \item Built a \textbf{unified benchmark} for large-scale traffic forecasting on the LargeST dataset, standardizing data processing and evaluation across classical ML, \textbf{spatiotemporal GNNs}, \textbf{State Space Models}, and \textbf{time-series foundation models} — using a \textbf{GIFT-Eval}-based pipeline that avoids constraining prediction horizons, letting every model compete on equal footing.
                \item Designed Mamba's selective-state SSM with a \textbf{parallel-scan algorithm} replacing sequential recurrence, improving accuracy by \textbf{23\%} and cutting training time by \textbf{35\%} versus the sequential baseline; further applied \textbf{in-training compression} to LRU, reducing training time by an additional \textbf{60\%} with minimal accuracy loss.


            \end{itemize}
        \end{onecolentry}

%--------------------------------------------------------------

\section*{Publications}
\begin{itemize}
  \item D.~Cao, \textbf{T.~Liao},Y.~Xu, N.~Ghasem, Q.~Guo, J.~Liu, M.~Weng, T.~Hirano, and Y.~Liu,
  ``Traffic Forecasting in Practice: Comparing Graph Models and Foundation Models
  for Deployment-Aware Traffic Forecasting,'' \textit{(under review / in submission)}, 2025.
\end{itemize}

\end{document}
